% titel: Blatt 5, A3: Radioaktiver Zerfall (C-14)
% Teil II: der durchgerechnete Weg.

\subsection*{Gegeben und gesucht}
$N(t)=N_0\mathrm e^{-\lambda t}$, Halbwertszeit
$T=\qty{5730}{\mathrm a}$, in einer Probe sind noch $q=53\,\%=0{,}53$ des
ursprünglichen C-14 vorhanden. Gesucht: $\lambda$ und das Alter $t$.

\subsection*{Ansatz und Gedankenschritt}
\emph{Zwei Schritte, zwei Unbekannte, zwei Gleichungen:} (1) $\lambda$ wird
aus der bekannten Halbwertszeit gewonnen, (2) das Alter aus dem gemessenen
Anteil. Im zweiten Schritt logarithmieren wir, weil $\ln$ die Umkehrung
von $e^{(\cdot)}$ ist und so das $t$ aus dem Exponenten \emph{holt}.

\subsection*{Zerfallskonstante $\lambda$}
\emph{Gedankenschritt:} Nach einer Halbwertszeit ist die Hälfte übrig.
Einsetzen in das Zerfallsgesetz und nach $\lambda$ auflösen:
\begin{equation*}
  \tfrac12 N_0=N_0\,\mathrm e^{-\lambda T}
  \;\Longrightarrow\;
  \ln\tfrac12=-\lambda T
  \;\Longrightarrow\;
  \lambda=\frac{\ln 2}{T}
  =\frac{\ln 2}{\qty{5730}{\mathrm a}}
  \approx 1{,}21\cdot10^{-4}\;\mathrm{a^{-1}}.
\end{equation*}
Richtiges Vorzeichen: $T$ hält $\lambda$ positiv, und die Menge nimmt ab.

\subsection*{Alter von Ötzi}
\emph{Gedankenschritt:} In $q=\mathrm e^{-\lambda t}$ isolieren wir $t$.
Setzen wir $\lambda=\ln 2/T$ ein, fällt das sofort ein; übrig bleibt ein
reiner Quotient aus Halbwertszeit, $\ln(1/q)$ und $\ln 2$.
\begin{equation*}
  0{,}53=\mathrm e^{-\lambda t}
  \;\Longrightarrow\;
  t=\frac{-\ln 0{,}53}{\lambda}
  =\frac{T\,\ln(1/0{,}53)}{\ln 2}.
\end{equation*}
Rechnung:
\begin{align*}
  t&=\frac{\qty{5730}{\mathrm a}\cdot\ln(1/0{,}53)}{\ln 2}
  =\frac{5730\cdot\ln(1{,}8868)}{0{,}6931} \\
  &\approx\frac{5730\cdot 0{,}6349}{0{,}6931}
  \approx \qty{5248}{\mathrm a}.
\end{align*}

\begin{ergebnis}
\begin{equation*}
  \lambda\approx 1{,}21\cdot10^{-4}\;\mathrm{a^{-1}},\qquad
  t\approx \qty{5248}{a}\ \text{(ca. }5{,}2\ \text{Millennien)} .
\end{equation*}
\end{ergebnis}

\begin{kontrolle}
\textbf{Grobenprüfung (die halbe Sache):} $53\,\%$ entspricht
$\ln(1/0{,}53)/\ln 2\approx0{,}916$ Halbwertszeiten. $0{,}916\cdot\qty{5730}{a}
\approx\qty{5250}{a}$ -- konsistent. \checkmark

\textbf{Plausibilität:} Ötzi wird auf ca.\ \qty{3300}{a}\ vor Christus
datiert; das liegt gut zu $5{,}2$ Millennien heute. Ein bekanntes,
glaubwürdiges Ergebnis. \checkmark

\textbf{Einheiten:} $\lambda$ in $\mathrm{a^{-1}}$, $t$ in $\unit{a}$ --
Produkt $\lambda t$ ist dimensionslos, wie es in einer Exponentialfunktion
sein muss. \checkmark
\end{kontrolle}

\begin{merksatz}
Aus der Halbwertszeit folgt $\lambda=\ln 2/T$. Aus dem Restanteil $q$:
$t=T\ln(1/q)/\ln 2$. $N_0$ kürzt sich heraus, nur der Anteil $q$ zählt;
mit $\ln$ den Exponenten auflösen. Einheiten ($\lambda$ und $t$) müssen
zusammenpassen.
\end{merksatz}
